\notepage{79}{Three-dimensional scattering}
The differential cross section relates the rate of scattering into a
small solid angle $\dd\Omega$ to the incident particle flux:
\[
 \frac{\dd\sigma}{\dd\Omega}\,\dd\Omega
 =\dd\sigma
 =\frac{\text{number scattered into }\dd\Omega\text{ per unit time}}
        {\text{number incident per unit time per unit area}}.
\]
To calculate it quantum mechanically, consider an incident plane wave
$\ee^{\ii kz}$ and an eigenfunction of $H=T+V$ of the form
\[
 \psi=\ee^{\ii kz}+\psi_{\rm sc}(r,\theta,\phi).
\]
The scattered wave is centered on the scattering center. Far outside
the potential, both terms satisfy the free Helmholtz equation at the
same energy:
\[
 (\nabla^2+k^2)\ee^{\ii kz}=0,\qquad
 (\nabla^2+k^2)\psi_{\rm sc}=0\quad(r\to\infty).
\]
The scattered wave can be expanded in spherical harmonics:
\[
 \psi_{\rm sc}\sim\sum_{\ell,m}
 \bigl[A_{\ell m}j_\ell(kr)+B_{\ell m}n_\ell(kr)\bigr]
 Y_\ell^m(\theta,\phi).
\]
Here $j_\ell$ and $n_\ell$ are spherical Bessel and Neumann functions.
Their large-$r$ limits are
\[
 j_\ell(kr)\sim\frac{\sin(kr-\ell\pi/2)}{kr},\qquad
 n_\ell(kr)\sim-\frac{\cos(kr-\ell\pi/2)}{kr}.
\]
Consequently, each radial combination has the asymptotic form
\[
 A_{\ell m}j_\ell(kr)+B_{\ell m}n_\ell(kr)
 \sim C_{\ell m}\frac{\ee^{\ii kr}}r
      +D_{\ell m}\frac{\ee^{-\ii kr}}r.
\]
Including the time factor gives
\[
 C_{\ell m}\frac{\ee^{\ii kr-\ii\omega t}}r
 +D_{\ell m}\frac{\ee^{-\ii kr-\ii\omega t}}r.
\]
The first term is outgoing, while the second is incoming. The physical
scattering setup chooses an outgoing scattered wave.
\editorial{The radial expansion coefficients are written with both
$\ell$ and $m$ indices for a general angular distribution. The source
abbreviates these coefficients by a single $\ell$ index. Its reading
reminder gives textbook page 348 without identifying an edition.}
